\chapter{Conclusiones} \label{chap:conclusiones}

En este trabajo hemos simulado el modelo de Ising 2D en red cuadrada con contorno periódico, comparando tres algoritmos de Montecarlo (Metropolis, Glauber y Wolff) para $L\in\{16,32,64,128\}$, con un doble objetivo: comprobar que los tres reproducen la física del equilibrio y comparar su comportamiento dinámico. En equilibrio los tres coinciden, como tienen que hacer al muestrear la misma distribución de Boltzmann. La estimación obtenida de $T_c$ por los cruces del cumulante de Binder es compatible con el valor de Onsager (con $U_4^*\approx 0{,}61$), y el colapso FSS de $\chi$ y $\xi_L/L$ con la búsqueda de los exponentes críticos confirma que los datos pertenecen a la clase de universalidad del modelo de Ising bidimensional. La única salvedad aparece cerca de $T_c$, en las magnitudes más sensibles a la autocorrelación, donde Metropolis y Glauber acusan la ralentización crítica, por lo que la convergencia limpia de $T_c$ con $L$ solo se ve en Wolff y el máximo de la susceptibilidad colapsada se desplaza a $x\approx 1$ en los algoritmos locales, frente a $x\approx 1{,}5$ en Wolff, que es el valor que tomamos como bueno.

La diferencia de fondo entre los algoritmos está en la dinámica, medida con el tiempo de autocorrelación integrado $\tau_{\text{int}}$. Metropolis y Glauber sufren ralentización crítica, con un máximo de $\tau_{\text{int}}$ que crece con $L$ en torno a $T_c$ de forma consistente con $z\approx 2$; dentro de esa clase, Glauber es algo más lento que Metropolis (un factor entre $1{,}5$ y $2$), porque Metropolis acepta los giros con probabilidad mayor o igual que la de Glauber para cada $\Delta E$ y descorrelaciona antes, lo cual es un resultado de validez general. Wolff, en cambio, prácticamente elimina la ralentización crítica: su $\tau_{\text{int}}$ es del orden de la unidad en todo el rango ($z\approx 0{,}017$), de modo que cerca de $T_c$ conseguir una muestra independiente cuesta varios órdenes de magnitud menos que con los algoritmos locales, una ventaja que se volvería decisiva con redes mayores.

Como continuación natural, sería interesante extender la comparación al Ising en tres dimensiones, donde no hay solución exacta y los exponentes solo se conocen numéricamente ahora que sabemos que el programa con el análisis de datos funciona de manera satisfactoria al comparar los resultados obtenidos con los resultados exactos del modelo de Ising bidimensional.

\newpage
\clearpage
\begin{otherlanguage}{english}
\noindent\textbf{\large Conclusions (English)}
\vspace{0.5em}

In this work we have simulated the 2D Ising model on a square lattice with periodic boundaries, comparing three Monte Carlo algorithms (Metropolis, Glauber and Wolff) for $L\in\{16,32,64,128\}$, with a twofold aim: to check that all three reproduce the equilibrium physics and to compare their dynamical behaviour. At equilibrium the three agree, as they must when sampling the same Boltzmann distribution. The estimate of $T_c$ obtained from the Binder cumulant crossings is compatible with Onsager's value (with $U_4^*\approx 0{,}61$), and the FSS collapse of $\chi$ and $\xi_L/L$ through the search for the critical exponents confirms that the data belong to the universality class of the two-dimensional Ising model. The only caveat appears near $T_c$, in the observables most sensitive to autocorrelation, where Metropolis and Glauber feel the critical slowing down, so that the clean convergence of $T_c$ with $L$ is seen only with Wolff and the peak of the collapsed susceptibility shifts to $x\approx 1$ for the local algorithms, against $x\approx 1{,}5$ for Wolff, which is the value we take as reliable.

The underlying difference between the algorithms lies in the dynamics, measured through the integrated autocorrelation time $\tau_{\text{int}}$. Metropolis and Glauber suffer critical slowing down, with a peak of $\tau_{\text{int}}$ that grows with $L$ around $T_c$ consistently with $z\approx 2$; within this class, Glauber is somewhat slower than Metropolis (by a factor between $1{,}5$ and $2$), because Metropolis accepts the flips with a probability greater than or equal to that of Glauber for every $\Delta E$ and decorrelates sooner, which is a general result. Wolff, by contrast, practically eliminates the critical slowing down: its $\tau_{\text{int}}$ is of unit order over the whole range ($z\approx 0{,}017$), so that near $T_c$ obtaining an independent sample costs several orders of magnitude less than with the local algorithms, an advantage that would become decisive for larger lattices.

As a natural continuation, it would be interesting to extend the comparison to the Ising model in three dimensions, where there is no exact solution and the exponents are known only numerically, now that we know that the program together with the data analysis works satisfactorily, as shown by comparing the results obtained against the exact results of the two-dimensional Ising model.
\end{otherlanguage}